\documentclass[11pt,a4paper]{article}
\usepackage[utf8]{inputenc}
\usepackage[french]{babel}
\usepackage[T1]{fontenc}
\usepackage{lmodern}
\usepackage[margin=1.5cm]{geometry}

\begin{document}

% En-tête
\begin{center}
    {\Huge \textbf{Lucas Moreau}} \\[0.2cm]
    {\Large \textbf{Analyste SOC}} \\[0.1cm]
    lucas.moreau@email.com \ | \ 06 10 22 14 69
\end{center}

\vspace{0.3cm}

% Résumé / Profil
\section*{Profil Professionnel}
\noindent
Ingénieur en cybersécurité passionné par la protection des infrastructures critiques et la traque des menaces avancées (Threat Hunting). Doté d'une solide expertise technique en test d'intrusion et en réponse à incident, je mets un point d'honneur à sensibiliser les équipes de développement aux enjeux du DevSecOps. Mon objectif est de garantir l'intégrité et la confidentialité des données sensibles face aux cyberattaques.

\vspace{0.3cm}

% Compétences
\section*{Compétences Techniques}
\noindent
\textbf{Sécurité Cloud (AWS/Azure)} $\bullet$ \textbf{Forensic} $\bullet$ \textbf{ISO 27001} $\bullet$ \textbf{Pentest} $\bullet$ \textbf{EDR / XDR} $\bullet$ \textbf{Python} $\bullet$ \textbf{Cryptographie} $\bullet$ \textbf{Splunk (SIEM)}

\vspace{0.3cm}

% Expériences
\section*{Expériences Professionnelles}
\noindent \textbf{Analyste SOC / Confirmé} --- \textit{Sogeti} \hfill \textbf{12/2023 à Aujourd'hui} \\
Implémentation d'une approche DevSecOps en intégrant des outils d'analyse statique et dynamique (SAST/DAST) directement dans les pipelines CI/CD. Reverse engineering de souches virales inédites pour extraire les indicateurs de compromission (IOC) et mettre à jour les bases de signatures. Réalisation de plus de 20 tests d'intrusion (pentests) sur des applications web et des infrastructures internes, avec rédaction de rapports d'audit détaillés. Investigation numérique (Forensic) suite à une compromission de serveur, identification de la porte dérobée (backdoor) et restauration de l'intégrité du système. \\[0.3cm]
\noindent \textbf{Analyste SOC} --- \textit{ANSSI} \hfill \textbf{01/2021 à 12/2023} \\
Accompagnement à la mise en conformité ISO 27001, incluant la rédaction des politiques de sécurité (PSSI) et la cartographie des risques. Gestion des incidents de sécurité de criticité majeure (P1) en coordination avec l'équipe de réponse à incident (CERT) et le comité de direction. Audit de configuration des équipements de sécurité (firewalls, proxys, WAF) pour s'assurer du strict respect des recommandations de l'ANSSI. Déploiement et configuration avancée d'une solution EDR (Endpoint Detection and Response) sur un parc de plus de 5 000 postes de travail. \\[0.3cm]
\noindent \textbf{Stage — Assistant Analyste SOC} --- \textit{Wavestone} \hfill \textbf{12/2020 à 01/2021} \\
Création de scripts d'automatisation en Python pour accélérer la collecte d'artefacts malveillants et l'analyse de logs réseau. Analyse en temps réel des alertes de sécurité remontées par le SIEM (Splunk), permettant la détection et le blocage d'attaques par ransomware. Animation de campagnes de sensibilisation à la cybersécurité (phishing simulé) ayant permis de réduire le taux de clics dangereux de 40\%. \\[0.3cm]


\vspace{0.3cm}

% Formations
\section*{Formation \& Diplômes}
\noindent \textbf{Diplôme d'Ingénieur en Réseaux et Sécurité} \hfill \textbf{2018 - 2021} \\
\textit{Télécom Paris} \\[0.3cm]
\noindent \textbf{Classes Préparatoires (CPGE) en Filière Scientifique} \hfill \textbf{2016 - 2018} \\
\textit{Lycée Scientifique} \\[0.3cm]


\end{document}